\documentclass{book}

\title{Analisi Matematica T-A}
\author{Michele Mazza}
\date{November 2023}


% Pacchetti di Codifica e Lingua
\usepackage[utf8]{inputenc}
\usepackage[T1]{fontenc}
\usepackage[italian]{babel}

% Pacchetti Grafici e Diagrammi
\usepackage{graphicx}
\usepackage{tikz}
\usepackage{tikz-3dplot}
\usetikzlibrary{angles,quotes,arrows.meta, positioning}
\usepackage{pgfplots}
\pgfplotsset{compat=1.18}

% Pacchetti Matematici
\usepackage{amsmath, amssymb, amsthm}

% Pacchetti per Tabelle e Unità di Misura
\usepackage{geometry}
\usepackage{siunitx}
\usepackage{booktabs}
\usepackage{tabularx}
\usepackage{colortbl}

% Pacchetti per Layout e Formattazione
\usepackage{fancyhdr}
\usepackage{float}
\usepackage{microtype}

% Altri Pacchetti Utili
\usepackage{mdframed}
\usepackage{animate}

% Pacchetto enumitem
\usepackage{enumitem}

% Definizioni dei Colori
\definecolor{lightblue}{HTML}{9cd8fc}
\definecolor{myblue}{HTML}{139FE9}

% Definizione del comando \norm
\newcommand{\norm}[1]{\left\lVert#1\right\rVert}

%\font\nullfont=cmr10 

% Ambienti Personalizzati
\newmdtheoremenv[
    backgroundcolor=orange!10,
    linecolor=myblue,
    linewidth=0pt,
    leftmargin=0pt,
    rightmargin=0pt,
    innerleftmargin=10pt,
    innerrightmargin=10pt,
    innertopmargin=10pt,
    innerbottommargin=10pt,
    nobreak
]{curiosity}{Curiosità}

\newmdtheoremenv[
    font=\normalfont\small,
    linewidth=0pt,
    innerleftmargin=10pt,
    innerrightmargin=10pt,
]{example}{Esempio}

\newmdtheoremenv[
  backgroundcolor=lightblue,
  linecolor=myblue,
  linewidth=1.5pt,
  innertopmargin=10pt,
  innerbottommargin=10pt,
  innerrightmargin=10pt,
  innerleftmargin=10pt,
  nobreak
]{theorem}{Teorema}

\newmdtheoremenv[
  backgroundcolor=gray!10,
  linewidth=0pt,
  innertopmargin=10pt,
  innerbottommargin=10pt,
  innerrightmargin=10pt,
  innerleftmargin=10pt,
  nobreak
]{definition}{Definizione}

\newmdtheoremenv[
  backgroundcolor=gray!10,
  linecolor=myblue,
  linewidth=1.5pt,
  innertopmargin=10pt,
  innerbottommargin=10pt,
  innerrightmargin=10pt,
  innerleftmargin=10pt,
  nobreak
]{proposition}{Proposizione}

% Configurazioni di fancyhdr
\pagestyle{fancy}
\fancyhf{} % Pulisce tutti i campi
\fancyhead[LE,RO]{\thepage} % Numero di pagina
\fancyhead[RE]{\nouppercase{\leftmark}} % Titolo del capitolo
\fancyhead[LO]{\nouppercase{\rightmark}} % Titolo della sezione
\renewcommand{\headrulewidth}{0.5pt}
\renewcommand{\footrulewidth}{0pt}

\geometry{a4paper, left=20mm, right=20mm, top=20mm, bottom=20mm}

% Pacchetto Hyperref (inserito per ultimo)
\usepackage{hyperref}

\begin{document}

\frontmatter
\maketitle
\tableofcontents
\mainmatter

\chapter{Proprietà globali delle funzioni continue}

\section{Teorema di esistenza degli zeri}
\begin{definition}
    \subsection{Zero di una funzione}
    Data una funzione reale \(f\), chiamiamo \textbf{zero} di \(f\) ogni punto \(x_0 \in dom f\) tale che \[f(x_0) = 0\]
\end{definition}

\begin{theorem}
    \subsection{T. di esistenza degli zeri}
    Sia \(\mathbf{f}\) una funzione \textbf{continua} in \(\mathbf{[a, b]}\) tale che \(\mathbf{f(a) \cdot f(b) < 0}\). Allora esiste \textbf{almeno un} punto \(\mathbf{x_0 \in (a, b)}\) tale che \(\mathbf{f(x_0) = 0}\).

    \noindent Ovvero se una funzione è continua in un intervallo chiuso e limitato e assume valori di segno opposto agli estremi, allora esiste almeno un punto in cui la funzione si annulla.
\end{theorem}

\begin{figure}[H]
    \centering
    \begin{tikzpicture}[scale = 1,
            declare function={
                    f(\x) = sin(deg(\x-2.5))*2;
                }]
        \begin{axis}
            [
                xlabel={\( x \)},
                ylabel={\( f(x) \)},
                xmin=-1, xmax=5,
                ymin=-3, ymax=3,
                grid=both,
                axis lines=middle,
                axis line style=->,
                grid style={line width=.1pt, draw=gray!10},
                major grid style={line width=.2pt,draw=gray!50},
                minor tick num=3,
                enlargelimits={abs=0.5},
                xlabel style={at={(ticklabel* cs:1)}, anchor=north west},
                ylabel style={at={(ticklabel* cs:1)}, anchor=south west},
                xticklabels={},
                yticklabels={}
                samples=100,
                unbounded coords=jump,
                legend pos=south east
            ]
            \draw [dashed, thick, color=gray] (axis cs: 0.25, {f(0.25)}) -- (axis cs: 0.25, 0);
            \draw [dashed, thick, color=gray] (axis cs: 0.25, {f(0.25)}) -- (axis cs: 0, {f(0.25)});
            \draw [dashed, thick, color=gray] (axis cs: 4.5, {f(4.5)}) -- (axis cs: 4.5, 0);
            \draw [dashed, thick, color=gray] (axis cs: 4.5, {f(4.5)}) -- (axis cs: 0, {f(4.5)});
            \addplot[domain=0.3:4.5, ultra thick, smooth, color=myblue]{f(x)};
            \filldraw[color=myblue] (axis cs: 0.3, {f(0.3)}) circle (1.5pt);
            \filldraw[color=myblue] (axis cs: 4.5, {f(4.5)}) circle (1.5pt);
            \filldraw[color=myblue] (axis cs: 2.5, {f(2.5)}) circle (1pt);
            \node[below right, color=myblue] at (axis cs:  2.4, {f(2.5)}) {\(x_0\)};
            \node[above, color=myblue] at (axis cs: 0.25, 0) {\(a\)};
            \node[below, color=myblue] at (axis cs: 4.5, 0) {\(b\)};
            \node[left, font=\small, color=myblue] at (axis cs: 0, {f(0.25)}) {\(f(a)\)};
            \node[left, font=\small, color=myblue] at (axis cs: 0, {f(4.5)}) {\(f(b)\)};
        \end{axis}
    \end{tikzpicture}
    \caption{Grafico di una funzione \(f\) che soddisfa le ipotesi del teorema.}
\end{figure}

\begin{theorem}
    [\textbf{Corollario 1 Unicità dello zero}]
    \noindent Se \(f\) è una funzione \textbf{continua e monotona} in \([a, b]\), e \(\mathbf{f(a) \cdot f(b) < 0}\), allora esiste un \textbf{unico} punto \(x_0 \in (a, b)\) tale che \(f(x_0) = 0\).
\end{theorem}

\begin{theorem}
    [\textbf{Corollario 2, Intersezione tra due funzioni}]
    Siano \(f, g\) due funzioni \textbf{continue} in \([a, b]\) tali che \(\mathbf{f(a) < g(a)}\) e \(\mathbf{f(b) > g(b)}\), allora esiste almeno un punto \(x_0 \in (a, b)\) tale che \(f(x_0) = g(x_0)\).
\end{theorem}
\begin{figure}[H]
    \centering
    \begin{tikzpicture}[scale = 1,
            declare function={
                    f(\x) = sin(deg(\x-2.5))*2+4;
                    g(\x) = sin(deg(\x-2.5))*3.4+4;
                }]
        \begin{axis}
            [
                xlabel={\( x \)},
                ylabel={\( f(x) \)},
                xmin=-1, xmax=5,
                ymin=-1, ymax=7,
                grid=both,
                axis lines=middle,
                axis line style=->,
                grid style={line width=.1pt, draw=gray!10},
                major grid style={line width=.2pt,draw=gray!50},
                minor tick num=3,
                enlargelimits={abs=0.5},
                xlabel style={at={(ticklabel* cs:1)}, anchor=north west},
                ylabel style={at={(ticklabel* cs:1)}, anchor=south west},
                xticklabels={},
                yticklabels={}
                samples=100,
                unbounded coords=jump,
                legend pos=south east
            ]
            \draw [dashed, thick, color=gray] (axis cs: 0.25, {f(0.25)}) -- (axis cs: 0.25, 0);
            \draw [dashed, thick, color=gray] (axis cs: 0.25, {f(0.25)}) -- (axis cs: 0, {f(0.25)});
            \draw [dashed, thick, color=gray] (axis cs: 4.5, {f(4.5)}) -- (axis cs: 4.5, 0);
            \draw [dashed, thick, color=gray] (axis cs: 4.5, {f(4.5)}) -- (axis cs: 0, {f(4.5)});
            \draw [dashed, thick, color=gray] (axis cs: 4.5, {g(4.5)}) -- (axis cs: 4.5, 0);
            \draw [dashed, thick, color=gray] (axis cs: 0.25, {g(0.25)}) -- (axis cs: 0, {g(0.25)});
            \draw [dashed, thick, color=gray] (axis cs: 4.5, {g(4.5)}) -- (axis cs: 0, {g(4.5)});
            \draw [dashed, thick, color=myblue] (axis cs: 2.5, {g(2.5)}) -- (axis cs: 2.5, 0);
            \addplot[domain=0.3:4.5, ultra thick, smooth, color=myblue]{f(x)};
            \addplot[domain=0.3:4.5, ultra thick, smooth, color=myblue]{g(x)};
            \filldraw[color=myblue] (axis cs: 0.3, {f(0.3)}) circle (1.5pt);
            \filldraw[color=myblue] (axis cs: 4.5, {f(4.5)}) circle (1.5pt);
            \filldraw[color=myblue] (axis cs: 0.3, {g(0.3)}) circle (1.5pt);
            \filldraw[color=myblue] (axis cs: 4.5, {g(4.5)}) circle (1.5pt);
            \filldraw[color=myblue] (axis cs: 2.5, {f(2.5)}) circle (2pt);
            \node[below right, color=myblue] at (axis cs:  2.4, 0) {\(x_0\)};
            \node[below] at (axis cs: 0.25, 0) {\(a\)};
            \node[below] at (axis cs: 4.5, 0) {\(b\)};
            \node[left, font=\small] at (axis cs: 0, {f(0.25)}) {\(f(a)\)};
            \node[left, font=\small] at (axis cs: 0, {f(4.5)}) {\(f(b)\)};
            \node[left, font=\small] at (axis cs: 0, {g(0.25)}) {\(g(a)\)};
            \node[left, font=\small] at (axis cs: 0, {g(4.5)}) {\(g(b)\)};
        \end{axis}
    \end{tikzpicture}
    \caption{Grafico di due funzioni \(f\) e \(g\) che soddisfano le ipotesi del teorema.}
\end{figure}

\section{Immagini di funzioni continue definite su intervalli}
In questo paragrafo studiamo l’immagine di una funzione continua definita su
un intervallo della retta reale.
\begin{theorem}
    \subsection{Teorema dei valori intermedi}
    Sia \(\mathbf{f}\) una funzione \textbf{continua} in \(\mathbf{[a, b]}\), chiuso e limitato allora \(\mathbf{f}\) assume tutti i valori compresi tra \(\mathbf{f(a)}\) e \(\mathbf{f(b)}\).
\end{theorem}
\begin{figure}[H]
    \centering
    \begin{tikzpicture}[scale = 1,
            declare function={
                    f(\x) = sin(deg(\x-2.5))*2+4;
                }]
        \begin{axis}
            [
                xlabel={\( x \)},
                ylabel={\( f(x) \)},
                xmin=-1, xmax=5,
                ymin=-1, ymax=7,
                grid=both,
                axis lines=middle,
                axis line style=->,
                grid style={line width=.1pt, draw=gray!10},
                major grid style={line width=.2pt,draw=gray!50},
                minor tick num=3,
                enlargelimits={abs=0.5},
                xlabel style={at={(ticklabel* cs:1)}, anchor=north west},
                ylabel style={at={(ticklabel* cs:1)}, anchor=south west},
                xticklabels={},
                yticklabels={}
                samples=100,
                unbounded coords=jump,
                legend pos=south east
            ]
            \draw [dashed, thick, color=gray] (axis cs: 0.25, {f(0.25)}) -- (axis cs: 0.25, 0);
            \draw [dashed, thick, color=gray] (axis cs: 0.25, {f(0.25)}) -- (axis cs: 0, {f(0.25)});
            \draw [dashed, thick, color=gray] (axis cs: 4.5, {f(4.5)}) -- (axis cs: 4.5, 0);
            \draw [dashed, thick, color=gray] (axis cs: 4.5, {f(4.5)}) -- (axis cs: 0, {f(4.5)});
            \addplot[domain=0.3:4.5, ultra thick, smooth, color=myblue]{f(x)};
            \filldraw[color=myblue] (axis cs: 0.3, {f(0.3)}) circle (1.5pt);
            \filldraw[color=myblue] (axis cs: 4.5, {f(4.5)}) circle (1.5pt);
            \filldraw[color=myblue] (axis cs: 2.5, {f(2.5)}) circle (1pt);
            \node[below] at (axis cs: 0.25, 0) {\(a\)};
            \node[below] at (axis cs: 4.5, 0) {\(b\)};
            \node[left, font=\small] at (axis cs: 0, {f(0.25)}) {\(f(a)\)};
            \node[left, font=\small] at (axis cs: 0, {f(4.5)}) {\(f(b)\)};
        \end{axis}
    \end{tikzpicture}
    \caption{Grafico di una funzione \(f\) che soddisfa le ipotesi del teorema.}
\end{figure}
\noindent\begin{minipage}{\textwidth}
    \begin{theorem}
        [\textbf{Corollario 1-Ponte da T. dei valori intermedi a T. di Weierstrass}]
        Se \( f \) è una funzione continua definita su un intervallo \( I \), allora l'insieme dei valori
        che \( f \) assume su \( I \) forma un intervallo che va da \( \inf_{x \in I} f(x) \) a
        \( \sup_{x \in I} f(x) \).
    \end{theorem}

    \begin{figure}[H]
        \centering
        \begin{tikzpicture}[scale = 1,
                declare function={
                        f(\x) = sin(deg(\x-2.5))*2+4;
                    }]
            \begin{axis}
                [
                    xlabel={\( x \)},
                    ylabel={\( f(x) \)},
                    xmin=-2, xmax=5,
                    ymin=-1, ymax=7,
                    grid=both,
                    axis lines=middle,
                    axis line style=->,
                    grid style={line width=.1pt, draw=gray!10},
                    major grid style={line width=.2pt,draw=gray!50},
                    minor tick num=3,
                    enlargelimits={abs=0.5},
                    xlabel style={at={(ticklabel* cs:1)}, anchor=north west},
                    ylabel style={at={(ticklabel* cs:1)}, anchor=south west},
                    xticklabels={},
                    yticklabels={}
                    samples=100,
                    unbounded coords=jump,
                    legend pos=south east
                ]
                \draw [dashed, thick, color=gray] (axis cs: 0.93, {f(0.93)}) -- (axis cs: 0.93, 0);
                \draw [dashed, thick, color=gray] (axis cs: 0.93, {f(0.93)}) -- (axis cs: 0, {f(0.93)});
                \draw [dashed, thick, color=gray] (axis cs: 4.07, {f(4.07)}) -- (axis cs: 4.07, 0);
                \draw [dashed, thick, color=gray] (axis cs: 4.07, {f(4.07)}) -- (axis cs: 0, {f(4.07)});
                \addplot[domain=0.3:4.5, ultra thick, smooth, color=myblue]{f(x)};
                \draw[ultra thick, color=myblue] (axis cs: 0.25,0) -- (axis cs: 4.5,0);
                \draw[ultra thick, color=myblue] (axis cs: 0,{f(0.93)}) -- (axis cs: 0,{f(4.07)});
                \filldraw[color=myblue] (axis cs: 0.93, {f(0.93)}) circle (1.5pt);
                \filldraw[color=myblue] (axis cs: 4.07, {f(4.07)}) circle (1.5pt);
                \node[below] at (axis cs: 0.25, 0) {\(a\)};
                \node[below] at (axis cs: 4.5, 0) {\(b\)};
                \node[below] at (axis cs: 0.93, 0) {\(x_m\)};
                \node[below] at (axis cs: 4.07, 0) {\(x_M\)};
                \node[left, font=\small] at (axis cs: 0, {f(0.93)}) {\(\inf_{x \in [a,b]} f\)};
                \node[left, font=\small] at (axis cs: 0, {f(4.07)}) {\(\sup_{x \in [a,b]} f\)};
            \end{axis}
        \end{tikzpicture}
        \caption{Grafico di una funzione \(f\) che soddisfa le ipotesi del teorema.}
    \end{figure}
\end{minipage}
\begin{theorem}
    \subsection{Teorema di Weierstrass}
    Sia \(\mathbf{f}\) una funzione \textbf{continua} in \(\mathbf{[a, b]}\), chiuso e limitato allora \(\mathbf{f}\) è limitata e ivi assume valori \textbf{minimo} e \textbf{massimo}.
    \begin{align*}
        m=min_{x \in [a,b]} f(x) \quad \text{e} \quad M=max_{x \in [a,b]} f(x)
    \end{align*}
    dunque
    \begin{align*}
        f[a,b]=[m,M]
    \end{align*}
\end{theorem}
\begin{figure}[H]
    \centering
    \begin{tikzpicture}[scale = 1,
            declare function={
                    f(\x) = sin(deg(\x-2.5))*2+4;
                }]
        \begin{axis}
            [
                xlabel={\( x \)},
                ylabel={\( f(x) \)},
                xmin=-1, xmax=5,
                ymin=-1, ymax=7,
                grid=both,
                axis lines=middle,
                axis line style=->,
                grid style={line width=.1pt, draw=gray!10},
                major grid style={line width=.2pt,draw=gray!50},
                minor tick num=3,
                enlargelimits={abs=0.5},
                xlabel style={at={(ticklabel* cs:1)}, anchor=north west},
                ylabel style={at={(ticklabel* cs:1)}, anchor=south west},
                xticklabels={},
                yticklabels={}
                samples=100,
                unbounded coords=jump,
                legend pos=south east
            ]
            \draw [dashed, thick, color=gray] (axis cs: 0.93, {f(0.93)}) -- (axis cs: 0.93, 0);
            \draw [dashed, thick, color=gray] (axis cs: 0.93, {f(0.93)}) -- (axis cs: 0, {f(0.93)});
            \draw [dashed, thick, color=gray] (axis cs: 4.07, {f(4.07)}) -- (axis cs: 4.07, 0);
            \draw [dashed, thick, color=gray] (axis cs: 4.07, {f(4.07)}) -- (axis cs: 0, {f(4.07)});
            \addplot[domain=0.3:4.5, ultra thick, smooth, color=myblue]{f(x)};
            \draw[ultra thick, color=myblue] (axis cs: 0.25,0) -- (axis cs: 4.5,0);
            \draw[ultra thick, color=myblue] (axis cs: 0,{f(0.93)}) -- (axis cs: 0,{f(4.07)});
            \filldraw[color=myblue] (axis cs: 0.93, {f(0.93)}) circle (1.5pt);
            \filldraw[color=myblue] (axis cs: 4.07, {f(4.07)}) circle (1.5pt);
            \node[below] at (axis cs: 0.25, 0) {\(a\)};
            \node[below] at (axis cs: 4.5, 0) {\(b\)};
            \node[below] at (axis cs: 0.93, 0) {\(x_m\)};
            \node[below] at (axis cs: 4.07, 0) {\(x_M\)};
            \node[left, font=\small] at (axis cs: 0, {f(0.93)}) {\(m\)};
            \node[left, font=\small] at (axis cs: 0, {f(4.07)}) {\(M\)};
        \end{axis}
    \end{tikzpicture}
    \caption{Grafico di una funzione \(f\) che soddisfa le ipotesi del teorema di Weierstrass.}
\end{figure}

\section{Invertibilità delle funzioni continue}
Nel capitolo 2 abbiamo visto che una funzione \(f\) è invertibile se è \textbf{biettiva}, ovvero se è \textbf{iniettiva} e \textbf{suriettiva}. In questo paragrafo studiamo le condizioni che permettono di affermare che una funzione continua è invertibile.
\begin{proposition}
    Sia \(f\) \textbf{continua} e \textbf{invertibile} su un intervallo \(I\), presi tre punti qualsiasi tali che \(x_1 < x_2 < x_3, \in I\), si verifica o una o l'altra delle seguenti relazioni:
    \begin{enumerate}
        \item \(f(x_1) < f(x_2) < f(x_3)\)
        \item \(f(x_1) > f(x_2) > f(x_3)\)
    \end{enumerate}
    Ciò deriva dalla proposizione che se \(f\) è invertibile allora è \textbf{monotona}.
\end{proposition}

\section{Funzioni Lipschitziane}
Le funzioni Lipschitziane sono funzioni che soddisfano una proprietà di \textbf{continuità} più forte rispetto a quella classica.
\begin{definition}
    \subsection{Funzione Lipschitziana}
    Sia \(f\) una funzione definita su un intervallo \(I\), diciamo che \(f\) è \textbf{Lipschitziana} su \(I\) se esiste una costante \(k > 0\) tale che
    \begin{align*}
        |f(x_1) - f(x_2)| \leq k|x_1 - x_2| \quad \forall x_1, x_2 \in I
    \end{align*}
    La costante \(k\) è detta \textbf{costante di Lipschitz} di \(f\).
\end{definition}
\textbf{Interpretazione}:
Una funzione Lipschitziana è una funzione che non cambia troppo rapidamente. Se sposti di un po' il punto su cui la stai calcolando, il suo valore non può aumentare o diminuire più di una certa quantità, che dipende dalla costante di Lipschitz.

\noindent\textbf{Applicazioni}:
Questo tipo di funzione è molto utile in matematica perché il suo comportamento è prevedibile e non presenta variazioni troppo drastiche. È spesso usata in campi come l'analisi numerica e le equazioni differenziali, dove la stabilità e la prevedibilità sono importanti.

\begin{proposition}
    Sia \(f\) una funzione \textbf{Lipschitziana} su un intervallo \(I\), allora \(f\) è \textbf{uniformemente continua} su \(I\).
\end{proposition}
\textbf{Uniformemente continua}:
Questo significa che la funzione non ha punti in cui cambia improvvisamente o drasticamente il suo comportamento. La continuità uniforme è una proprietà più forte della semplice continuità, poiché implica un tipo di regolarità e prevedibilità in tutto l'intervallo su cui la funzione è definita.

\begin{theorem}
    \subsection{Teorema di Heine-Cantor}
    Sia \(f\) una funzione \textbf{continua} su un intervallo \textbf{chiuso e limitato} \(I\), allora \(f\) è \textbf{uniformemente continua} su \(I\).
\end{theorem}
\end{document}